% from ega3-1-fr.tex
\oldpage[III]{121}
un $\mathcal S$-Module gradué \emph{de type fini}. En effet, la question est
locale sur $\mathfrak X$, et on peut donc supposer pour la traiter que
$\mathfrak X=\operatorname{Spf}(B)$, où $B$ est un anneau adique noethérien,
et $\mathcal F=N^\Delta$, où $N$ est un $B$-module de type fini
(I, 10.10.5)~; on a en outre $X_0=\operatorname{Spec}(B_0)$ où
$B_0=B/\mathfrak J B$, et les $\mathcal O_{X_0}$-Modules quasi-cohérents
$\mathcal S$ et $\mathcal M$ sont respectivement égaux à $\widetilde{S'}$ et
$\widetilde{M'}$, où
\[
S'=\bigoplus_{k\geq0}
((\mathfrak J^k/\mathfrak J^{k+1})\otimes_{A_0}B_0),\qquad
M'=\bigoplus_{k\geq0}
((\mathfrak J^k/\mathfrak J^{k+1})\otimes_{A_0}N_0),
\]
avec $N_0=N/\mathfrak JN$~; on a donc évidemment
$M'=S'\otimes_{B_0}N_0$, et comme $N_0$ est un $B_0$-module de type fini,
$M'$ est un $S'$-module de type fini, d'où notre assertion (I, 1.3.13).

Comme le morphisme $f_0:X_0\to Y_0$ est \emph{propre} par hypothèse, on peut
appliquer (\hyperref[III.3.3.2-fr]{3.3.2}) à $\mathcal S$, $\mathcal M$ et au
morphisme $f_0$~; compte tenu de (\hyperref[III.1.4.11-fr]{1.4.11}), on en
conclut que pour \emph{tout $n\geq0$},
$\bigoplus_{k\geq0}H^n(X_0,\mathcal M_k)$ est un $S$-module gradué \emph{de
type fini}. Cela prouve que la condition $(\mathrm F_n)$ de (0, 13.7.7) est
vérifiée pour \emph{tout $n\geq0$}, lorsqu'on considère le système projectif
strict $(\mathcal F/\mathfrak J^k\mathcal F)_{k\geq0}$ de faisceaux de
groupes abéliens sur $X_0$, munis chacun de sa structure naturelle de
«~$A$-module filtré~». On peut donc appliquer (0, 13.7.7), qui prouve que~:
\begin{enumerate}
\item[$1^\circ$] Le système projectif
$(H^n(\mathfrak X,\mathcal F_k))_{k\geq0}$ vérifie la condition (ML).
\item[$2^\circ$] Si
$H'^{,n}=\varprojlim_kH^n(\mathfrak X,\mathcal F_k)$, $H'^{,n}$ est un
$A$-module de type fini.
\item[$3^\circ$] La filtration définie sur $H'^{,n}$ par les noyaux des
homomorphismes canoniques $H'^{,n}\to H^n(\mathfrak X,\mathcal F_k)$ est
$\mathfrak J$-bonne.
\end{enumerate}

Notons d'autre part que si on pose
$X_k=(\mathfrak X,\mathcal O_{\mathfrak X}/\mathcal K^{k+1})$,
$\mathcal F_k$ est un $\mathcal O_{X_k}$-Module cohérent (I, 10.11.3) et si
$U$ est un ouvert affine dans $X_0$, $U$ est aussi un ouvert affine dans
chacun des $X_k$ (I, 5.1.9), donc $H^n(U,\mathcal F_k)=0$ pour tout $n>0$ et
tout $k$ (\hyperref[III.1.3.1-fr]{1.3.1}) et
$H^0(U,\mathcal F_k)\to H^0(U,\mathcal F_h)$ est surjectif pour $h\leq k$
(I, 1.3.9). On est donc dans les conditions de (0, 13.3.2) et l'application
de (0, 13.3.1) prouve que $H'^{,n}$ s'identifie canoniquement à
$H^n(\mathfrak X,\varprojlim_k\mathcal F_k)=H^n(\mathfrak X,\mathcal F)$~;
cela achève de prouver (\hyperref[III.3.4.4-fr]{3.4.4}).

\begin{env}[3.4.5]
\label{III.3.4.5-fr}
Revenons maintenant à la démonstration de (\hyperref[III.3.4.2-fr]{3.4.2})
et (\hyperref[III.3.4.3-fr]{3.4.3}). Prouvons d'abord ces propositions dans
le cas $\mathfrak Y=\operatorname{Spf}(A)$ envisagé dans
(\hyperref[III.3.4.4-fr]{3.4.4})~; pour cela, pour tout $g\in A$, appliquons
(\hyperref[III.3.4.4-fr]{3.4.4}) au schéma formel affine noethérien induit
sur l'ouvert $\mathfrak Y_g=\mathfrak D(g)$ de $\mathfrak Y$, qui est égal à
$\operatorname{Spf}(A_{\{g\}})$, et au préschéma formel induit par
$\mathfrak X$ sur $f^{-1}(\mathfrak Y_g)$~; notons que $\mathfrak Y_g$ est
aussi un ouvert affine dans le préschéma
$Y_k=(\mathfrak Y,\mathcal O_{\mathfrak Y}/(\mathfrak J^\Delta)^{k+1})$,
et comme $\mathcal F_k$ est un $\mathcal O_{X_k}$-Module cohérent, on a
\[
H^n(f^{-1}(\mathfrak Y_g),\mathcal F_k)
=\Gamma(\mathfrak Y_g,R^nf_*(\mathcal F_k))
\]
pour tout $k\geq0$ en vertu de
(\hyperref[III.1.4.11-fr]{1.4.11}). L'homomorphisme canonique
\[
H^n(f^{-1}(\mathfrak Y_g),\mathcal F)
\longrightarrow\varprojlim_k\Gamma(\mathfrak Y_g,R^nf_*(\mathcal F_k))
\]
est un isomorphisme~; mais on a (0\textsubscript{I}, 3.2.6)
\[
\varprojlim_k\Gamma(\mathfrak Y_g,R^nf_*(\mathcal F_k))
=\Gamma(\mathfrak Y_g,\varprojlim_kR^nf_*(\mathcal F_k))
\]

